\section{Fermi Data Analysis}
\label{sec:fermi-data-analysis}


In this section, we describe our procedure for measuring the spectrum and intensity of the gamma-ray emission from each of the Milky Way's 157 globular clusters~\cite{harris2010newcatalogglobularclusters}.
We employ the python package, Fermipy~\cite{wood2017fermipyopensourcepythonpackage}, as provided by the Fermi collaboration, to process the data collected by the Fermi-LAT. 
We work with HEALPix maps of order 11, corresponding to approximately $0.03^{\circ}$ bins, and utilize 15.8 years of data.\footnote{MET range:  239557414 - 733276805} We consider photons in the range of 100 MeV to 100 GeV, divided into 5 logarithmic bins per decade. We adopt the standard event selection for gamma-ray point sources, using the \verb|SOURCE| event class and \verb|FRONT| + \verb|BACK| event types. We also employ a zenith angle cut of $90^{\circ}$ and a rocking angle cut of $52^{\circ}$ in order to reduce the contribution from the Earth's limb. We require data quality=1 and only consider data that was taken while in survey mode. We do not use data that was collected while passing through the South Atlantic Anomaly. These and other aspects of our analysis are summarized in Table~\ref{tab:settings}.

\begin{table}[t]
\centering
\begin{tabular}{ | l| l|  }
\hline
HEALPix map order& 11  \\ 
Weeks & $9-826$  \\ 
Emin & 100 MeV \\
Emax & 100 GeV \\
Instrument Response Functions (IRFs) & \verb|P8R3_SOURCE_V3| \\
EVCLASS & 128 (Source)  \\ 
EVTYPE & 3 (Front + Back)  \\
ZMAX & $90^{\circ}$ \\
Rocking angle cut & $52^{\circ}$ \\
Galactic diffuse emission template & \verb|gll_iem_v07| \\
Isotropic background template & \verb|iso_P8R3_SOURCE_V3_v1| \\
Source catalog & \verb|gll_psc_v35| \\ 
\hline
\end{tabular}
\caption{{\tt Fermi Science Tools} settings used in this analysis. }
\label{tab:settings}
\end{table}

For each globular cluster, we study a $14^{\circ} \times 14^{\circ}$ region-of-interest (ROI) centered around the nominal position of the source. We then fit the energy spectrum of the globular cluster, adopting a parameterization in the form of a power law with an exponential cut-off:
\begin{equation}
    \phi(E) = N_0 \, E^{-\alpha} \, e^{-E/E_c}. 
\label{eq:energy-spectrum}
\end{equation}
%
In fitting these sources, we follow the standard procedure used in constructing Fermi-LAT source catalogs. 
We include in our fit any sources in the 4FGL-DR4 catalog that are located within in ROI~\cite{ballet2024fermilargeareatelescope}, excluding any globular cluster, pulsar, or unassociated gamma-ray source that is located within $0.2^\circ$ of the fiducial position of the globular cluster in question. We then add a new point source in the position of the globular cluster, as specified in Ref.~\cite{harris2010newcatalogglobularclusters}, and reoptimize the model of the known sources within the ROI.
Following standard procedure, we run the source localization pipeline, readjust the centroid of the globular cluster, and look for new point sources with a power-law spectrum. 
We then move to the characterization of the globular cluster, performing a 3D likelihood scan in the ($\alpha$, $E_{c}$, $N_0$) parameter space. 
%
During this procedure, we model the Galactic diffuse emission and isotropic background contributions according to the templates provided by the Fermi collaboration (see Table \ref{tab:settings}). We allow their normalization to float, as well as the energy spectrum normalization of any other resolved source within 3 degrees.
%
This allows us to identify, for each value of the integrated flux, the combination of these parameters that yields the greatest likelihood. 

In Table~\ref{tab:high_TS}, which we report in Appendix \ref{sec:app-tables}, we list each of the 56 globular clusters for which our analysis identified a statistically significant test statistic, \mbox{TS > 25}, along with the best-fit values for $\alpha$, $E_c$, and the total flux (integrated between 0.1 and 100 GeV). Note that our previous study, which utilized 7 years of Fermi data, detected only 25 globular clusters at a level of \mbox{TS > 25}~\cite{Hooper:2016rap}. We have compared our results to those reported in the 4FGL catalog~\cite{ballet2024fermilargeareatelescope} and find them to be in overall good agreement. The only significant differences are found in the cases of NGC 5139, Palomar 6, Terzan 5, and NGC 6218, for which our analysis yields integrated fluxes which are larger than those given in the 4FGL catalog by $2-3\sigma$. 

\noindent
This appears to be due to the spectral parameterization we have adopted, as these four sources are fit in the 4FGL to a log-parabola or power-law parameterization, rather than that of a power-law with an exponential cutoff.

In addition to the 56 sources with TS > 25 listed in Table~\ref{tab:high_TS}, we found an additional 7 sources which yielded TS > 25, but that had their highest TS values in bins below 0.4 GeV. Due to the large point spread function and related uncertainties that are encountered at these low energies, we do not consider these sources (which are listed in Table~\ref{tab:low_peak_TS}) to be unambiguous detections, and only report upper limits on their flux.

Finally, in Table~\ref{tab:low_TS}, we report upper limits on the gamma-ray fluxes from each of the remaining globular clusters considered in our analysis, each of which yielded TS < 25. 

